\origpage{233}{GRAPH MINORS IV}
\hspace*{12pt}(3) \textit{There is an infinite $M$-stable set $X \subseteq V(R_1 \cup R_2 \cup \dots)$ such that $|X \cap V(T_i)| \le 1$ for each $i \ge 1$, and the set $B$ of heads of all edges of $R_1 \cup R_2 \cup \dots$ with tails in $X$ is $M$-rich.}

We claim that $X$ satisfies the theorem. For let $C$ be the set of heads of all edges of $T$ with tails in $X$, and let $z \in C$. Then $z$ is the head of some edge $f$ of some $T_i$ with tail in $X$. If $\phi_i(f) = 0$ then $f \in F_i$ and so $z \in A$. If $\phi_i(f) > 0$ then $f$ is an edge of some $R_j$, since its tail is in some $V(R_j)$, and so $z \in B$. Thus $C \subseteq A \cup B$. But $A$ is $M$-rich from (2), and $B$ is $M$-rich from (3), and so $C$ is $M$-rich. Thus $X$ satisfies the theorem. $\square$

To illustrate the use of (2.2) we show that it (indeed, (2.1)) implies Kruskal's theorem (1.2). Let $T_1, T_2, \dots$ be a countable sequence of disjoint trees, and let $M$ be the infinite graph with $V(M) = V(T_1 \cup T_2 \cup \dots)$ in which for $i' > i \ge 1$, $u \in V(T_i)$ is $M$-adjacent to $v \in V(T_{i'})$ if (with notation as in (2.1)) $T^v$ topologically contains $T^u$. (We define topological containment for directed graphs in the natural way; however, we do not demand that root be mapped to root.) Now the conclusion of (2.1) cannot be satisfied, by Higman's ``finite sequence'' theorem ((8.3), later); and so its hypothesis that $\{o(T_1), o(T_2), \dots\}$ is $M$-stable is false. This proves Kruskal's theorem.

One may derive similarly a theorem of Friedman [12] using (2.2) in place of (2.1). By adjusting the definition of $M$ accordingly, one may also derive the versions of these theorems when the vertices are labelled from a well-quasi-order.

\sectionhead{3. Hypergraphs and Tree-Decompositions}

A hypergraph $G$ consists of a finite set $V(G)$ of vertices, a finite set $E(G)$ of edges with $E(G) \cap V(G) = \emptyset$, and an incidence relation, a subset of $V(G) \times E(G)$. The vertices incident with an edge $e$ are its \textit{ends}, and the set of ends of $e$ in $G$ is denoted by $V(e)$ (or $V_G(e)$ in cases of ambiguity). If $G$, $G'$ are hypergraphs and $V(G') \subseteq V(G)$ and $E(G') \subseteq E(G)$, and $V_G(e) = V_{G'}(e)$ for every $e \in E(G')$, then we call $G'$ a \textit{subhypergraph} of $G$ and write $G' \subseteq G$. If $G_1, G_2 \subseteq G$ we define $G_1 \cap G_2$, $G_1 \cup G_2$ in the natural way. A \textit{separation} of $G$ is a pair $(G_1, G_2)$ of subhypergraphs of $G$ with $G_1 \cup G_2 = G$ and $E(G_1 \cap G_2) = \emptyset$; its \textit{order} is $|V(G_1 \cap G_2)|$.

A march in a set $V$ is a finite sequence of distinct elements of $V$. If $\pi$ is the march $v_1, \dots, v_k$ we denote $\{v_1, \dots, v_k\}$ by $\bar{\pi}$. A \textit{rooted hypergraph} $G = (G^-, \pi(G))$ consists of a hypergraph $G^-$ and a march $\pi(G)$ in $V(G^-)$. We define $V(G) = V(G^-)$, $E(G) = E(G^-)$. A separation of a rooted hypergraph $G$ is a pair $(G_1, G_2)$ of rooted hypergraphs such that $(G_1^-, G_2^-)$ is a separation of $G^-$, and $\pi(G_2) = \pi(G)$, and $\bar{\pi}(G_1) = V(G_1^- \cap G_2^-)$. Its order is
